\chapter{Conclusioni}
\label{cap:conclusioni}

\section{Sintesi dei risultati}
\label{sec:sintesi}

% DA COMPLETARE: da scrivere per ultimo, quando i risultati della simulazione e della
% valutazione economica saranno disponibili. Di seguito i risultati gia' acquisiti, che
% andranno integrati.

Il lavoro svolto finora ha prodotto un risultato metodologico: \`e possibile ricostruire il
prezzo del Mercato del Giorno Prima a partire dai soli dati pubblici delle offerte, con
un'accuratezza sufficiente a usarlo come base per il calcolo di prezzi controfattuali. Sulle
744 aste orarie di gennaio 2025 il prezzo ricostruito coincide esattamente con quello
ufficiale nel 52{,}3\% dei casi, \`e entro 1~\euromwh{} nell'83{,}7\% ed entro 5~\euromwh{}
nel 99{,}5\%, con mediana dello scarto nulla, deviazione standard di 1{,}06~\euromwh{} e uno
scarto massimo di 6{,}61~\euromwh{} su tutto il mese.

Arrivare a questo risultato ha richiesto cinque elementi, nessuno dei quali evidente
all'inizio: la selezione delle offerte che hanno effettivamente partecipato all'asta;
l'inclusione delle offerte presentate a granularit\`a diversa da quella prevalente;
l'uso della quantit\`a rettificata al posto di quella offerta; il trattamento delle offerte a
blocchi come indivisibili; e soprattutto la reintroduzione dello scambio netto con l'esterno
del perimetro zonale, senza il quale il prezzo ricostruito sbaglia di circa 100~\euromwh{}.

Lo strumento che ne risulta permette di misurare la grandezza su cui poggia l'intera tesi:
di quanto si sposta il prezzo di equilibrio quando si inietta potenza nel mercato. \`E la
stessa grandezza che determina il ricavo di un accumulo e che ne limita la scala, perch\'e
una batteria che carica alza il prezzo a cui compra e che scarica abbassa quello a cui vende.

Su questa base \`e stata definita la misura centrale del lavoro, l'\textbf{erosione di
profitto}~\eqref{eq:erosione}: la quota di margine che l'accumulo distrugge da s\'e, ottenuta
confrontando lo stesso piano valorizzato ai prezzi senza accumulo e ai prezzi che l'accumulo
stesso genera. La soglia $K^*$ a cui il quantile prudenziale dell'erosione supera un livello
dichiarato caratterizza il passaggio da price taker a price maker come fenomeno aleatorio,
con il suo intervallo di confidenza.

% DA COMPLETARE: sintesi dei risultati quantitativi definitivi, da riscrivere quando il
% campione sara' esteso ad almeno un anno. Il risultato preliminare su gennaio 2025 colloca
% K* a 136 MW con intervallo 83-179, un valore basso rispetto ai 15-25 GW scambiati per
% asta: la ragione e' che l'accumulo opera nelle ore estreme, dove la curva e' piu' ripida,
% e non nell'ora media.

\paragraph{Il confronto con la letteratura.} Il risultato va letto accanto a quello di
\textcite{alonsoperez2026storage}, che studiano la stessa transizione sul mercato spagnolo e
individuano soglie di saturazione attorno ai 15 e ai 32~GWh. I due numeri non sono
confrontabili direttamente — mercati, anni e perimetri sono diversi, e le loro soglie sono
espresse in energia su scala nazionale mentre queste sono in potenza su una singola zona — e
il confronto va inteso in termini qualitativi.

La differenza che conta \`e un'altra ed \`e metodologica. La loro stima \`e \textbf{puntuale},
perch\'e il modello sottostante \`e deterministico; quella proposta qui \`e una
\textbf{distribuzione}, con un intervallo di confidenza che nel campione di gennaio 2025 va da
83 a 179~MW, cio\`e largo pi\`u del doppio del suo estremo inferiore. Riportare un solo numero
avrebbe suggerito una precisione che i dati non hanno, e avrebbe nascosto proprio
l'informazione che serve a un investitore, che deve dimensionare sul caso avverso ragionevole
e non su quello tipico.

\textcite{veenstra2025profitability} giungono, sul mercato olandese, a una conclusione pi\`u
netta: il valore dell'arbitraggio \textbf{si satura} al crescere della capacit\`a installata, al
punto che anche ipotizzando volatilit\`a elevata dei prezzi e una riduzione dei costi di
investimento del 60\% resterebbe poco spazio per nuova capacit\`a profittevole. \`E un risultato
della letteratura, non di questo lavoro, ma \`e quello con cui i risultati italiani dovranno
dialogare pi\`u direttamente, perch\'e riguarda la stessa grandezza qui misurata dall'erosione
di profitto.

Un elemento di contatto \`e gi\`a visibile: anche in zona NORD l'erosione cresce pi\`u che
proporzionalmente con la capacit\`a, fino a rendere il profitto negativo
(Sezione~\ref{sec:autodanno}), il che \`e a tutti gli effetti una forma di saturazione. Se il
confronto quantitativo confermi o contraddica la loro conclusione resta per\`o questione aperta.

% DA COMPLETARE: confronto quantitativo con la conclusione di saturazione, da svolgere quando
% la curva di erosione sara' disponibile su un campione annuale.

\section{Limiti del modello e sviluppi futuri}
\label{sec:limiti}

Si elencano i limiti gi\`a identificati e misurati, distinguendo quelli di cui \`e noto il
costo da quelli ancora aperti.

\paragraph{Perimetro zonale e scambio calibrato.} La zona NORD \`e modellata senza vincoli
espliciti di transito, e l'energia scambiata con l'esterno \`e reintrodotta come blocco
calibrato sull'esito osservato dell'asta. Ne consegue che la simulazione dell'accumulo assume
flussi di scambio insensibili alla variazione di prezzo che l'accumulo stesso induce:
l'assunzione \`e tanto pi\`u forte quanto pi\`u il periodo \`e congestionato. Il superamento
naturale \`e un modello multi-zonale con i limiti di transito espliciti.

\paragraph{Offerte a blocchi.} Il vincolo di accettazione tutto-o-niente \`e rappresentato
con un'euristica iterativa che non garantisce l'ottimo del problema intero sottostante. Il
confronto con la soluzione osservata dell'asta mostra che l'approssimazione \`e buona: sulle
aste orarie la eguaglia, su quelle a quarto d'ora ne recupera l'86\%.

\paragraph{Residuo non spiegato nelle aste a quarto d'ora.} Restano 133~MW per asta, pari
allo 0{,}74\% del volume scambiato, non attribuiti ad alcuna causa nota. Due ipotesi sono
state formulate e scartate sulla base dei dati — l'accettazione parziale delle offerte e
l'esistenza di vincoli fra quarti d'ora consecutivi — e la ricerca resta aperta.

\paragraph{Reazione di secondo ordine non catturata.} \`E il limite pi\`u rilevante per
l'interpretazione dei risultati. Gli altri operatori del mercato non modificano le proprie
offerte in risposta all'ingresso dell'accumulo: le curve d'asta restano quelle osservate
storicamente. Nella realt\`a un produttore che vedesse comprimersi il prezzo serale
riformulerebbe le proprie offerte, e cos\`i farebbe un consumatore flessibile. L'erosione
misurata \`e quindi l'effetto \emph{di primo ordine}, e non \`e possibile stabilire con questo
impianto se la reazione degli altri operatori la attenuerebbe o l'amplificherebbe.

\paragraph{Aggiustamento intra-periodo non catturato.} Le batterie non ricontrattano la
propria posizione dopo la chiusura del Mercato del Giorno Prima, mentre nella realt\`a
opererebbero anche sui mercati infragiornalieri, dove potrebbero recuperare parte del margine
eroso. La stima dell'erosione sul solo MGP \`e quindi conservativa rispetto alla redditivit\`a
complessiva, ma non necessariamente rispetto alla soglia.

\paragraph{Previsione perfetta.} Il profilo di carica e scarica \`e ottimizzato conoscendo i
prezzi di tutti i periodi della giornata. Il profitto che ne risulta \`e un limite superiore,
non un risultato conseguibile da un operatore reale. Poich\'e per\`o l'erosione \`e un rapporto
fra due profitti calcolati sullo stesso piano, un piano subottimo ridurrebbe entrambi i
termini e l'effetto sulla soglia \`e di secondo ordine.

\paragraph{Stima in coda del quantile.} La soglia \`e definita su un quantile della
distribuzione dell'erosione fra i giorni. Anche restando all'80\textsuperscript{o} o al
90\textsuperscript{o} percentile — e si \`e evitato deliberatamente di spingersi oltre — la
stima poggia su un numero limitato di giornate, e gli intervalli di confidenza bootstrap
riportati insieme a $K^*$ servono a rendere esplicita questa fragilit\`a.

\paragraph{Flotta omogenea.} La capacit\`a aggregata \`e trattata come un'unica batteria
equivalente. L'equivalenza vale finch\'e le batterie hanno gli stessi vincoli, lineari, e
ottimizzano sullo stesso segnale di prezzo; si rompe se le durate o gli stati iniziali sono
diversi, perch\'e in quel caso la somma dei profili ottimi non coincide con il profilo ottimo
della somma.

\paragraph{Esistenza dell'equilibrio nella variante a operatore unico.} Nella formulazione
alternativa, in cui un operatore unico riottimizza tenendo conto del proprio effetto, profilo
e prezzi si determinano a vicenda e il punto fisso pu\`o non esistere. \`E una delle ragioni
per cui quella formulazione non \`e stata adottata, ma resta implementata come termine di
confronto: il divario fra i due scenari misura quanto valga, per l'accumulo, internalizzare
il proprio effetto sul prezzo.

\paragraph{Mercati non considerati.} L'analisi riguarda il solo Mercato del Giorno Prima.
Un sistema di accumulo reale opera anche sui Mercati Infragiornalieri e sul Mercato per il
Servizio di Dispacciamento, dove i ricavi possono essere significativi: la redditivit\`a
stimata in questo lavoro va quindi letta come una componente, non come il totale.

% DA COMPLETARE: aggiungere i limiti che emergeranno dalla simulazione e dalla valutazione
% economica, in particolare quelli relativi al modello di degrado e al trattamento
% dell'incertezza sui prezzi.
